\documentclass[10pt,a4paper]{article}
\usepackage[margin=18mm]{geometry}
\usepackage{amsmath,amssymb,mathtools}
\usepackage{hyperref}
\hypersetup{hidelinks,hypertexnames=false}
\allowdisplaybreaks
\sloppy
\setlength{\emergencystretch}{2em}
\title{Odd denominators in power-sum expansions\\Stanley supplementary exercise 21}
\author{Statement: Richard P. Stanley\\Author proof: Darij Grinberg}
\date{Source answer 403672, edited 30 September 2021}
\begin{document}
\maketitle
\section*{Statement (Stanley, supplementary exercise 21)}
Let \(f\in\Lambda_{\mathbb Z}\cap\mathbb Q[p_1,p_3,p_5,\ldots]\), and write
\(f=\sum_\lambda a_\lambda p_\lambda\).
Show that when the rational number \(a_\lambda\) is written in lowest terms,
its denominator is odd.

\paragraph{Source correspondence.}
Stanley's solutions explicitly cite Darij Grinberg's answer to MathOverflow question 403153.
The author's Theorem 1 below treats an arbitrary prime \(q\);
the original exercise uses its \(q=2\) case.
All supporting results supplied in that answer are retained here.

\section*{Author answer (Darij Grinberg)}


Here is a proof of the generalization suggested by Richard Stanley in the
comments and even of a more general result (with "odd" replaced by "not
divisible by a given prime $q$"). It is completely different from the argument
I sketched in the comments, and is completely elementary (using no Macdonald
polynomials). Unfortunately, it is also somewhat awkward and way too long
(much of it devoted to a fight with notations).




For any commutative ring $R$, we let $\Lambda_{R}$ be the ring of symmetric
functions over $R$; this is a commutative $R$-algebra. Let
$\operatorname{Par}$ be the set of all partitions. We shall use the standard
notation $p_{\lambda}$ for the power-sum symmetric function indexed by a
partition $\lambda$.




Fix a prime $q$. Let $\mathbb{Z}_{\left(  q\right)  }$ denote the ring of all
rational numbers that can be written in the form $\dfrac{a}{b}$ for two
integers $a$ and $b$ such that $b$ is coprime to $q$. These numbers are known
as \emph{$q$-integers}. Obviously, $\mathbb{Z}_{\left(  q\right)  }$ is a subring
of $\mathbb{Q}$, so that $\Lambda_{\mathbb{Z}_{\left(  q\right)  }}$ is a
subring of $\Lambda_{\mathbb{Q}}$.




Now, Richard Stanley's generalization (generalized a bit further) claims:



\begin{quote}



\textbf{Theorem 1.} We have
\begin{align*}
\Lambda_{\mathbb{Z}_{\left(  q\right)  }}\cap\mathbb{Q}\left[  p_{i}
\ \mid\ i\not \equiv 0\mod q\right]  =\mathbb{Z}_{\left(
q\right)  }\left[  p_{i}\ \mid\ i\not \equiv 0\mod q\right]  .
\end{align*}



\end{quote}



The proof requires some preparations, which include showing some results of
independent interest.




First, we introduce a few more classical notations from symmetric function
theory: For any partition $\lambda$ and any $i\geq1$, we let $m_{i}\left(
\lambda\right)  $ denote the multiplicity of $i$ in $\lambda$ (that is, the
number of parts of $\lambda$ that equal $i$). For any partition $\lambda$, we
define the positive integer
\begin{align*}
z_{\lambda}:=\prod_{i=1}^{\infty}\left(  \left(  m_{i}\left(  \lambda\right)
\right)  !\cdot i^{m_{i}\left(  \lambda\right)  }\right)  .
\end{align*}




Let $V$ be the $\Lambda_{\mathbb{Z}}$-subalgebra of $\Lambda_{\mathbb{Q}}$
generated by the fractions $\dfrac{p_{i}}{i^{k}}$ for all positive integers
$i$ and all nonnegative integers $k$. In other words, let
\begin{align*}
V=\Lambda_{\mathbb{Z}}\left[  \dfrac{p_{i}}{i^{k}}\ \mid\ i>0\text{ and }
k\geq0\right]  \subseteq\Lambda_{\mathbb{Q}}.
\end{align*}




Now, we claim the following:



\begin{quote}



\textbf{Theorem 2.} We have
\begin{align*}
z_{\lambda}^{-1}p_{\lambda}\in V
\end{align*}
for each partition $\lambda$.



\end{quote}



To prove this, we need a simple arithmetic lemma:



\begin{quote}



\textbf{Lemma 3.} Let $c$ and $d$ be two integers with $d\neq0$. Then, there exist
some integers $a$ and $b$ and some nonnegative integer $i$ such that
$c^i =ad+bc^{i+1}$.



\end{quote}



\emph{Proof of Lemma 3.} The ring $\mathbb{Z}/d\mathbb{Z}$ is finite (since
$d\neq0$). For any integer $m$, we let $\overline{m}\in\mathbb{Z}/d\mathbb{Z}$
denote the residue class of $m$ in this ring. The infinitely many residue
classes $\overline{c^{0}},\overline{c^{1}},\overline{c^{2}},\ldots$ all belong
to the finite ring $\mathbb{Z}/d\mathbb{Z}$, and thus cannot all be distinct
(by the pigeonhole principle). In other words, there exist two nonnegative
integers $i$ and $j$ satisfying $i<j$ and $\overline{c^i }=\overline{c^j }$.
Consider these $i$ and $j$. We have $\overline{c^i }=\overline{c^j }$; in
other words, $c^i \equiv c^j \mod d$. Hence, $d\mid c^i 
-c^j $. In other words, $c^i -c^j =ad$ for some integer $a$. Consider this
$a$. However, recall that $i<j$. Thus, $j=i+w$ for some positive integer $w$.
Consider this $w$. The integer $w-1$ is nonnegative (since $w$ is positive);
thus, $c^{w-1}$ is an integer. Hence, we can define an integer $b$ by
$b=c^{w-1}$. From $j = i+w = \left(w-1\right) + \left(i+1\right)$, we obtain $c^j = c^{\left(w-1\right) + \left(i+1\right)} = \underbrace{c^{w-1}}_{= b} c^{i+1} = bc^{i+1}$.
Now, from $c^i -c^j =ad$, we obtain
\begin{align*}
c^i   & =ad+\underbrace{c^j }_{= bc^{i+1}} =ad+bc^{i+1}.
\end{align*}
This proves Lemma 3. $\blacksquare$




\emph{Proof of Theorem 2.} We shall use the notation $\ell\left(  \lambda\right)  $
for the \emph{length} of a partition $\lambda$ (that is, the number of all parts of
$\lambda$). For instance, $\ell\left(  \left(  5,2,2\right)  \right)  =3$. We
shall also use the notation $\left\vert \lambda\right\vert $ for the \emph{size} of
a partition $\lambda$ (that is, the sum of all parts of $\lambda$). We shall
prove Theorem 2 by strong induction on $\left\vert \lambda\right\vert
+\ell\left(  \lambda\right)  $. Thus, we fix some $N\in\mathbb{N}$, and we
assume (as the induction hypothesis) that Theorem 2 holds for all $\lambda$
with $\left\vert \lambda\right\vert +\ell\left(  \lambda\right)  <N$. We now
must prove Theorem 2 for all $\lambda$ with $\left\vert \lambda\right\vert
+\ell\left(  \lambda\right)  =N$.




So let $\lambda$ be a partition satisfying $\left\vert \lambda\right\vert
+\ell\left(  \lambda\right)  =N$. We must prove that $z_{\lambda}
^{-1}p_{\lambda}\in V$.




If $\lambda=\varnothing$, then this is obvious (since $z_{\lambda}
^{-1}p_{\lambda}=1$ in this case). Thus, for the rest of this induction step,
we WLOG assume that $\lambda\neq\varnothing$.




We are in one of the following three cases:




\emph{Case 1:} All parts of $\lambda$ are equal to $1$.




\emph{Case 2:} All parts of $\lambda$ are equal, but not equal to $1$.




\emph{Case 3:} Not all parts of $\lambda$ are equal.




Let us consider Case 1 first. In this case, all parts of $\lambda$ are equal
to $1$. In other words, $\lambda=\underbrace{\left(  1,1,\ldots,1\right)
}_{v\text{ entries}}$ for some positive integer $v$ (since $\lambda
\neq\varnothing$). Consider this $v$. Thus, $p_{\lambda}=p_{1}^v $ and
$\left\vert \lambda\right\vert =v$ and $\ell\left(  \lambda\right)  =v$.
Let $\operatorname{Par}_{v}$ denote the set of all
partitions of $v$. Thus, $\lambda\in\operatorname{Par}_{v}$. Now, let $h_v \in \Lambda_{\mathbb{Z}}$ denote the $v$-th complete homogeneous symmetric function. A
well-known formula (e.g., (2.5.17) in \href{https://arxiv.org/abs/1409.8356v7}{Grinberg/Reiner, arXiv:1409.8356v7})
yields
\begin{align*}
h_{v}=\sum_{\mu\in\operatorname{Par}_{v}}z_{\mu}^{-1}p_{\mu
}=z_{\lambda}^{-1}p_{\lambda}+\sum_{\substack{\mu\in\operatorname{Par}_{v};
\\\mu\neq\lambda}}z_{\mu}^{-1}p_{\mu},
\end{align*}
so that
\begin{equation}
z_{\lambda}^{-1}p_{\lambda}=h_v -\sum_{\substack{\mu\in
\operatorname{Par}_{v};\\\mu\neq\lambda}}z_{\mu}^{-1}p_{\mu}.
\label{darij1.pf.t2.c1.2}
\tag{1}
\end{equation}




However, the only partition of $v$ that has length $\geq v$ is the partition
$\underbrace{\left(  1,1,\ldots,1\right)  }_{v\text{ entries}}=\lambda$. Thus,
if $\mu$ is a partition of $v$ distinct from $\lambda$, then $\mu$ has length
$<v$. In other words, if $\mu\in\operatorname{Par}_{v}$ satisfies
$\mu\neq\lambda$, then $\ell\left(  \mu\right)  <v$. Hence, if $\mu
\in\operatorname{Par}_{v}$ satisfies $\mu\neq\lambda$, then
\begin{align*}
\underbrace{\left\vert \mu\right\vert }_{=v=\left\vert \lambda\right\vert
}+\underbrace{\ell\left(  \mu\right)  }_{<v=\ell\left(  \lambda\right)
}<\left\vert \lambda\right\vert +\ell\left(  \lambda\right)  =N
\end{align*}
and therefore $z_{\mu}^{-1}p_{\mu}\in V$ (by our induction hypothesis, applied
to $\mu$ instead of $\lambda$). Thus, \eqref{darij1.pf.t2.c1.2} becomes
\begin{align*}
z_{\lambda}^{-1}p_{\lambda}=\underbrace{h_v}_{\in\Lambda_{\mathbb{Z}
}\subseteq V}-\sum_{\substack{\mu\in\operatorname{Par}_{v};\\\mu
\neq\lambda}}\underbrace{z_{\mu}^{-1}p_{\mu}}_{\in V}\in V-\sum_{\substack{\mu
\in\operatorname{Par}_{v};\\\mu\neq\lambda}}V\subseteq V.
\end{align*}
Hence, $z_{\lambda}^{-1}p_{\lambda}\in V$ has been proved in Case 1.




Let us next consider Case 2. In this case, all parts of $\lambda$ are equal,
but not equal to $1$. In other words, $\lambda=\underbrace{\left(
n,n,\ldots,n\right)  }_{v\text{ entries}}$ for some positive integers $n\neq1$
and $v$ (since $\lambda\neq\varnothing$). Consider these $n$ and $v$. Thus,
$p_{\lambda}=p_n^v $ and $z_{\lambda}=v!\cdot n^v $ and $\left\vert
\lambda\right\vert =nv$ and $\ell\left(  \lambda\right)  =v$. From $n\neq1$,
we obtain $n>1$ (since $n$ is a positive integer). Thus, $nv > v$ (since $v > 0$). In other words, $v < nv$.




From $p_{\lambda}=p_n^v $ and $z_{\lambda}=v!\cdot n^v $, we obtain
\begin{equation}
z_{\lambda}^{-1}p_{\lambda}=\left(  v!\cdot n^v \right)  ^{-1}p_n
^v =\dfrac{p_n^v }{v!\cdot n^v }.
\label{darij1.pf.t2.c2.1}
\tag{2}
\end{equation}




We must prove that $z_{\lambda}^{-1}p_{\lambda}\in V$. In other words, we must
prove that $\dfrac{p_n^v }{v!\cdot n^v }\in V$ (since $z_{\lambda}
^{-1}p_{\lambda}=\dfrac{p_n^v }{v!\cdot n^v }$).




We shall now use Lemma 3 to decompose $\dfrac{p_n^v }{v!\cdot n^v }$ into
a sum of partial fractions -- one with a denominator of $v!$ and another with
a power of $n$ in the denominator. We will then prove that both of these
fractions belong to $V$.




Indeed, Lemma 3 (applied to $c=n^v $ and $d=v!$) yields that there exist some
integers $a$ and $b$ and some nonnegative integer $i$ such that $\left(
n^v \right)  ^i =a\cdot v!+b\left(  n^v \right)  ^{i+1}$. Consider these
$a$, $b$ and $i$. Multiplying both sides of the equality $\left(
n^v \right)  ^i =a\cdot v!+b\left(  n^v \right)  ^{i+1}$ by $\dfrac
{p_n^v }{\left(  n^v \right)  ^{i+1}\cdot v!}$, we obtain
\begin{align}
\dfrac{p_n^v }{v!\cdot n^v }
&=\left(a\cdot v!+b\left(  n^v \right)
^{i+1}\right) \cdot \dfrac{p_n^v }{\left(  n^v \right)  ^{i+1}\cdot v!} \\
&=a\cdot\dfrac{p_n^v }{\left(
n^v \right)  ^{i+1}}+b\cdot\dfrac{p_n^v }{v!}
.
\label{darij1.pf.t2.c2.parfrac}
\tag{3}
\end{align}
Thus, in order to prove that $\dfrac{p_n^v }{v!\cdot n^v }\in V$, it
suffices to show that $\dfrac{p_n^v }{\left(  n^v \right)  ^{i+1}}\in V$
and $\dfrac{p_n^v }{v!}\in V$ (because $a$ and $b$ are integers, and $V$ is
a ring).




The first of these two claims is easy: We have $\dfrac{p_n}{n^{i+1}}\in V$
(since $\dfrac{p_n}{n^{i+1}}$ is one of the designated generators of the
$\Lambda_{\mathbb{Z}}$-algebra $V$). Hence, $\left(  \dfrac{p_n}{n^{i+1}
}\right)  ^v \in V$ (since $V$ is a ring). In other words, $\dfrac{p_n^v 
}{\left(  n^v \right)  ^{i+1}}\in V$ (since $\left(  \dfrac{p_n}{n^{i+1}
}\right)  ^v =\dfrac{p_n^v }{\left(  n^v \right)  ^{i+1}}$).




It remains to prove that $\dfrac{p_n^v }{v!}\in V$. To do this, we will
need the \emph{Frobenius endomorphism} $\mathbf{f}_n$. It is defined as follows:
For any commutative ring $R$, we let
\begin{align*}
\mathbf{f}_n:\Lambda_{R}\rightarrow\Lambda_{R}
\end{align*}
be the $R$-algebra homomorphism that sends each symmetric function $f$ to
$f\left(  x_{1}^{n},x_{2}^{n},x_{3}^{n},\ldots\right)  $ (where we regard $f$
as a symmetric formal power series in countably many indeterminates $x_{1},x_{2}
,x_{3},\ldots$). This homomorphism $\mathbf{f}_n$ is called the $n$-th
\emph{Frobenius endomorphism} and is functorial in $R$ (that is, it commutes with
the morphisms $\Lambda_{R}\rightarrow\Lambda_{S}$ induced by ring
homomorphisms $R\rightarrow S$). If you like to think in terms of plethysm,
$\mathbf{f}_n$ can be described as sending each $f\in\Lambda_{R}$ to the
plethysm $f\left[  p_n\right]  $.




The functoriality of $\mathbf{f}_n$ in $R$ entails that the $\mathbf{f}_n$
defined for $R=\mathbb{Z}$ is a restriction of the $\mathbf{f}_n$ defined
for $R=\mathbb{Q}$. Thus, we can safely denote both of these maps by
$\mathbf{f}_n$ without risking confusion. They both are ring homomorphisms
(since they are $R$-algebra homomorphisms for appropriate $R$). Of course,
$\mathbf{f}_n\left(  \Lambda_{\mathbb{Z}}\right)  \subseteq\Lambda
_{\mathbb{Z}}$ (since the $\mathbf{f}_n$ defined for $R=\mathbb{Z}$ is a
restriction of the $\mathbf{f}_n$ defined for $R=\mathbb{Q}$).




It is easy to see that $\mathbf{f}_n\left(  p_{i}\right)  =p_{in}$ for each
$i>0$. Hence, for each positive integer $i$ and each nonnegative integer $k$,
we have
\begin{align*}
\mathbf{f}_n\left(  \dfrac{p_{i}}{i^{k}}\right)  =\dfrac{p_{in}}{i^{k}
}=n^{k}\cdot\dfrac{p_{in}}{\left(  in\right)  ^{k}}\in V
\end{align*}
(since $\dfrac{p_{in}}{\left(  in\right)  ^{k}}$ is one of the designated
generators of the $\Lambda_{\mathbb{Z}}$-algebra $V$). Therefore,
\begin{align*}
\Lambda_{\mathbb{Z}}\left[  \mathbf{f}_n\left(  \dfrac{p_{i}}{i^{k}}\right)
\ \mid\ i>0\text{ and }k\geq0\right]  \subseteq V
\end{align*}
(since $V$ is a $\Lambda_{\mathbb{Z}}$-algebra).




Now, from $V=\Lambda_{\mathbb{Z}}\left[  \dfrac{p_{i}}{i^{k}}\ \mid\ i>0\text{
and }k\geq0\right]  $, we obtain
\begin{align*}
\mathbf{f}_n\left(  V\right)   &  =\mathbf{f}_n\left(  \Lambda
_{\mathbb{Z}}\left[  \dfrac{p_{i}}{i^{k}}\ \mid\ i>0\text{ and }k\geq0\right]
\right)  \\
&  =\underbrace{\left(  \mathbf{f}_n\left(  \Lambda_{\mathbb{Z}}\right)
\right)  }_{\subseteq\Lambda_{\mathbb{Z}}}\left[  \mathbf{f}_n\left(
\dfrac{p_{i}}{i^{k}}\right)  \ \mid\ i>0\text{ and }k\geq0\right]  \\
&  \qquad\left(  \text{since }\mathbf{f}_n\text{ is a ring homomorphism}
\right)  \\
&  \subseteq\Lambda_{\mathbb{Z}}\left[  \mathbf{f}_n\left(  \dfrac{p_{i}
}{i^{k}}\right)  \ \mid\ i>0\text{ and }k\geq0\right]  \subseteq V.
\end{align*}




Let $\mu$ be the partition $\underbrace{\left(  1,1,\ldots,1\right)
}_{v\text{ entries}}$. Then, $p_{\mu}=p_{1}^v $ and $z_{\mu}=v!\cdot
\underbrace{1^v }_{=1}=v!$ and $\left\vert \mu\right\vert =v$ and
$\ell\left(  \mu\right)  =v$. Hence,
\begin{align*}
\underbrace{\left\vert \mu\right\vert }_{=v}+\underbrace{\ell\left(
\mu\right)  }_{=v}  & =\underbrace{v}_{< nv} + v\\
& <\underbrace{nv}_{=\left\vert \lambda\right\vert }+\underbrace{v}
_{=\ell\left(  \lambda\right)  }=\left\vert \lambda\right\vert +\ell\left(
\lambda\right)  =N.
\end{align*}
Hence, $z_{\mu}^{-1}p_{\mu}\in V$ (by our induction hypothesis, applied to
$\mu$ instead of $\lambda$). In view of $z_{\mu}=v!$ and $p_{\mu}=p_{1}^v $,
this rewrites as $v!^{-1}\cdot p_{1}^v \in V$. Hence,
\begin{equation}
\mathbf{f}_n\left(  v!^{-1}\cdot p_{1}^v \right)  \in \mathbf{f}
_n\left(  V\right)  \subseteq V.
\label{darij1.pf.t2.c2.5}
\tag{4}
\end{equation}
However, since $\mathbf{f}_n$ is a $\mathbb{Q}$-algebra homomorphism, we
have
\begin{align*}
\mathbf{f}_n\left(  v!^{-1}\cdot p_{1}^v \right)  =v!^{-1}\cdot\left(
\mathbf{f}_n\left(  p_{1}\right)  \right)  ^v =\dfrac{\left(
\mathbf{f}_n\left(  p_{1}\right)  \right)  ^v }{v!}=\dfrac{p_n^v }{v!}
\end{align*}
(since $\mathbf{f}_n\left(  p_{1}\right)  =p_n$). Thus,
\eqref{darij1.pf.t2.c2.5} rewrites as $\dfrac{p_n^v }{v!}\in V$. Hence,
\eqref{darij1.pf.t2.c2.parfrac} becomes
\begin{align*}
\dfrac{p_n^v }{v!\cdot n^v }=a\cdot\underbrace{\dfrac{p_n^v }{\left(
n^v \right)  ^{i+1}}}_{\in V}+b\cdot\underbrace{\dfrac{p_n^v }{v!}}_{\in
V}\in V
\end{align*}
(since $V$ is a ring and since $a$ and $b$ are integers). In view of
\eqref{darij1.pf.t2.c2.1}, this rewrites as $z_{\lambda}^{-1}p_{\lambda}\in
V$. Hence, $z_{\lambda}^{-1}p_{\lambda}\in V$ has been proved in Case 2.




Let us finally consider Case 3. In this case, not all parts of $\lambda$ are equal.




We need another notation: If $\alpha=\left(  \alpha_{1},\alpha_{2}
,\ldots,\alpha_{m}\right)  $ and $\beta=\left(  \beta_{1},\beta_{2}
,\ldots,\beta_n\right)  $ are two partitions, then $\alpha\sqcup\beta$ shall
denote the partition obtained by sorting the tuple $\left(  \alpha_{1}
,\alpha_{2},\ldots,\alpha_{m},\beta_{1},\beta_{2},\ldots,\beta_n\right)  $
in weakly decreasing order. For instance, $\left(  3,2,2\right)  \sqcup\left(
5,3,2\right)  =\left(  5,3,3,2,2,2\right)  $.




It is easy to see that if $\alpha$ and $\beta$ are two partitions that have no
part in common, then
\begin{equation}
z_{\alpha\sqcup\beta}=z_{\alpha}z_{\beta}.
\label{darij1.pf.t2.c3.zaub}
\tag{5}
\end{equation}
Moreover, if $\alpha$ and $\beta$ are any two partitions, then
\begin{equation}
p_{\alpha\sqcup\beta}=p_{\alpha}p_{\beta}.
\label{darij1.pf.t2.c3.paub}
\tag{6}
\end{equation}




Now, recall that not all parts of $\lambda$ are equal. Hence, we can write
$\lambda$ in the form $\lambda=\alpha\sqcup\beta$ where $\alpha$ and $\beta$
are two nonempty partitions that have no part in common. (Indeed, we can
define $\alpha$ and $\beta$ by choosing an arbitrary part $i$ of $\lambda$,
then letting $\alpha$ be the partition consisting of all parts of $\lambda$
equal to $i$, while $\beta$ is the partition consisting of all remaining parts
of $\lambda$.) Consider these $\alpha$ and $\beta$. It is easy to see that
$\left\vert \alpha\right\vert <\left\vert \alpha\sqcup\beta\right\vert $
(since $\beta$ is nonempty) and $\ell\left(  \alpha\right)  <\ell\left(
\alpha\sqcup\beta\right)  $ (for the same reason). Since $\alpha\sqcup
\beta=\lambda$, these two inequalities rewrite as $\left\vert \alpha
\right\vert <\left\vert \lambda\right\vert $ and $\ell\left(  \alpha\right)
<\ell\left(  \lambda\right)  $. Adding these two inequalities together, we
obtain
\begin{align*}
\left\vert \alpha\right\vert + \ell\left(  \alpha\right)
 <\left\vert \lambda\right\vert +\ell\left(  \lambda\right)  =N.
\end{align*}
Hence, $z_{\alpha}^{-1}p_{\alpha}\in V$ (by our induction hypothesis, applied
to $\alpha$ instead of $\lambda$). Similarly, $z_{\beta}^{-1}p_{\beta}\in V$.
However, from \eqref{darij1.pf.t2.c3.zaub} and \eqref{darij1.pf.t2.c3.paub},
we obtain
\begin{align*}
z_{\alpha\sqcup\beta}^{-1}p_{\alpha\sqcup\beta}=\left(  z_{\alpha}z_{\beta
}\right)  ^{-1}p_{\alpha}p_{\beta}=\underbrace{z_{\alpha}^{-1}p_{\alpha}}_{\in
V}\cdot\underbrace{z_{\beta}^{-1}p_{\beta}}_{\in V}\in V
\end{align*}
(since $V$ is a ring). In view of $\alpha\sqcup\beta=\lambda$, this rewrites
as $z_{\lambda}^{-1}p_{\lambda}\in V$. Hence, $z_{\lambda}^{-1}p_{\lambda}\in
V$ has been proved in Case 3.




We have now proved $z_{\lambda}^{-1}p_{\lambda}\in V$ in all three cases 1, 2
and 3. Thus, the induction step is complete.




Thus, Theorem 2 is proved by induction. $\blacksquare$




We are not quite ready to prove Theorem 1 yet. We first need some more notations.




We let $\operatorname{QPar}$ be the set of all partitions that have no part
divisible by $q$. (If $q=2$, these are precisely the partitions into odd parts.)




We let $J$ be the ideal of the ring $\Lambda_{\mathbb{Q}}$ generated by the
$p_{i}$ with $i\equiv0\mod q$. In other words, $J=\sum\limits_{i=1}^{\infty}p_{iq}\Lambda_{\mathbb{Q}}$. Recall that the family $\left(
p_{\lambda}\right)  _{\lambda\in\operatorname{Par}}$ is a basis of the
$\mathbb{Q}$-vector space $\Lambda_{\mathbb{Q}}$. Thus, the $\mathbb{Q}
$-vector subspace $J$ of $\Lambda_{\mathbb{Q}}$ has basis $\left(  p_{\lambda
}\right)  _{\lambda\in\operatorname{Par}\setminus\operatorname{QPar}}$
(because multiplying any $p_{\mu}$ by a $p_{iq}$ yields a $p_{\lambda}$ with
$\lambda\in\operatorname{Par}\setminus\operatorname{QPar}$, and conversely,
any $p_{\lambda}$ with $\lambda \in \operatorname{Par}\setminus\operatorname{QPar}$ can
be obtained in such a way).




A well-known fact (or easy exercise) in abstract algebra says the following:



\begin{quote}



\textbf{Lemma 4.} Let $B$ be a subring of a ring $A$. Let $I$ be a (two-sided)
ideal of $A$. Then, $B+I$ is a subring of $A$.



\end{quote}



Applying Lemma 4 to $A=\Lambda_{\mathbb{Q}}$, $B=\Lambda_{\mathbb{Z}_{\left(
q\right)  }}$ and $I=J$, we conclude that $\Lambda_{\mathbb{Z}_{\left(
q\right)  }}+J$ is a subring of $\Lambda_{\mathbb{Q}}$. We denote this subring
$\Lambda_{\mathbb{Z}_{\left(  q\right)  }}+J$ by $W$. We note that $W$ is
furthermore a $\mathbb{Z}_{\left(  q\right)  }$-subalgebra of $\Lambda
_{\mathbb{Q}}$ (since $W$ is a subring of $\Lambda_{\mathbb{Q}}$ and is
preserved under scaling by $\mathbb{Z}_{\left(  q\right)  }$).




Next, we observe:



\begin{quote}



\textbf{Proposition 5.} We have $V\subseteq W$.



\end{quote}



\emph{Proof of Proposition 5.} We have
\begin{align}
\Lambda
_{\mathbb{Z}}\subseteq \Lambda
_{\mathbb{Z}_{\left(  q\right)  }}\subseteq \Lambda_{\mathbb{Z}_{\left(
q\right)  }}+J=W .
\end{align}
Now, $W$ is a commutative ring (since it is a subring of $\Lambda_{\mathbb{Q}}$) and contains
$\Lambda_{\mathbb{Z}}$ as a subring (since $\Lambda
_{\mathbb{Z}}\subseteq W$). Hence, $W$ is a $\Lambda_{\mathbb{Z}}$-algebra. Thus, in order to prove that $V\subseteq W$, it suffices to show
that $\dfrac{p_{i}}{i^{k}}\in W$ for each positive integer $i$ and each
nonnegative integer $k$ (by the definition of $V$).




So let us show this. Fix a positive integer $i$ and a nonnegative integer $k$. We must prove that $\dfrac{p_{i}}{i^{k}}\in W$. If $i\equiv0\mod q$,
then this follows from the obvious fact that $\dfrac{p_{i}}{i^{k}}\in
J\subseteq\Lambda_{\mathbb{Z}_{\left(  q\right)  }}+J=W$. Thus, we WLOG assume that $i\not \equiv 0\mod q$. Hence, $i$ is coprime to $q$. Hence, $\dfrac{1}{i}\in \mathbb{Z}_{\left(  q\right)  }$, so that $\dfrac{1}{i^{k}}\in\mathbb{Z}_{\left(  q\right)  }\subseteq W$. Now, $\dfrac{p_{i}}{i^{k}} =\underbrace{\left(  \dfrac{1}{i}\right)  ^{k}}_{\in W} \underbrace{p_{i}}_{\in \Lambda_{\mathbb{Z}} \subseteq W}\in W$ (since $W$ is a ring). This completes our proof of Proposition 5.
$\blacksquare$




At last, we can prove Theorem 1.




\emph{Proof of Theorem 1.} It is clear that $\mathbb{Z}_{\left(  q\right)  }\left[ p_{i}\ \mid\ i\not \equiv 0\mod q\right]  \subseteq
\Lambda_{\mathbb{Z}_{\left(  q\right)  }}\cap\mathbb{Q}\left[  p_{i}
\ \mid\ i\not \equiv 0\mod q\right]  $. Hence, it suffices to
prove the reverse inclusion, i.e., to prove that
\begin{align*}
\Lambda_{\mathbb{Z}_{\left(  q\right)  }}\cap\mathbb{Q}\left[  p_{i}
\ \mid\ i\not \equiv 0\mod q\right]  \subseteq \mathbb{Z}_{\left(  q\right)  }\left[ p_{i}\ \mid\ i\not \equiv 0\mod q\right] .
\end{align*}
Thus, we fix an arbitrary $f\in\Lambda_{\mathbb{Z}_{\left(  q\right)  }}
\cap\mathbb{Q}\left[  p_{i}\ \mid\ i\not \equiv 0\mod q\right]
$. We must prove that $f\in \mathbb{Z}_{\left(  q\right)  }\left[ p_{i}\ \mid\ i\not \equiv 0\mod q\right]$.




We have $f\in\Lambda_{\mathbb{Z}_{\left(  q\right)  }}\cap\mathbb{Q}\left[
p_{i}\ \mid\ i\not \equiv 0\mod q\right]  \subseteq
\mathbb{Q}\left[  p_{i}\ \mid\ i\not \equiv 0\mod q\right]  $.
Hence, $f$ is a $\mathbb{Q}$-linear combination of the family $\left(
p_{\lambda}\right)  _{\lambda\in\operatorname{QPar}}$ (since this family
$\left(  p_{\lambda}\right)  _{\lambda\in\operatorname{QPar}}$ is a basis of
the $\mathbb{Q}$-vector space $\mathbb{Q}\left[  p_{i}\ \mid\ i\not \equiv
0\mod q\right]  $). In other words, we can write $f$ in the form
\begin{equation}
f=\sum_{\lambda\in\operatorname{QPar}}c_{\lambda}p_{\lambda}
\label{darij1.pf.t1.f=}
\tag{7}
\end{equation}
for some $c_{\lambda}\in\mathbb{Q}$. Consider these $c_{\lambda}$. We shall prove that they all belong to $\mathbb{Z}_{\left(q\right)}$.




We let $\left\langle \cdot,\cdot\right\rangle $ denote the Hall inner product
on $\Lambda_{\mathbb{Q}}$. This is a $\mathbb{Q}$-bilinear form sending
$\Lambda_{\mathbb{Q}}\times\Lambda_{\mathbb{Q}}$ to $\mathbb{Q}$ and sending
$\Lambda_{\mathbb{Z}_{\left(  q\right)  }}\times\Lambda_{\mathbb{Z}_{\left(
q\right)  }}$ to $\mathbb{Z}_{\left(  q\right)  }$ (since its restriction to
$\Lambda_{\mathbb{Z}_{\left(  q\right)  }}\times\Lambda_{\mathbb{Z}_{\left(
q\right)  }}$ is the Hall inner product on $\Lambda_{\mathbb{Z}_{\left(
q\right)  }}$).




It is well-known (see, e.g., Corollary 2.5.17(b) in \href{https://arxiv.org/abs/1409.8356v7}{Grinberg/Reiner,
arXiv:1409.8356v7}) that the families $\left(  p_{\lambda}\right)
_{\lambda\in\operatorname{Par}}$ and $\left(  z_{\lambda}^{-1}p_{\lambda
}\right)  _{\lambda\in\operatorname{Par}}$ are dual bases of $\Lambda
_{\mathbb{Q}}$ with respect to the Hall inner product. Hence,
\begin{equation}
\left\langle p_{\lambda},z_{\mu}^{-1}p_{\mu}\right\rangle =\delta_{\lambda
,\mu}
\label{darij1.pf.t1.dualbases}
\tag{8}
\end{equation}
for any $\lambda\in\operatorname{Par}$ and $\mu\in\operatorname{Par}$ (where
the $\delta_{\lambda,\mu}$ is a Kronecker delta). In other words,
\begin{equation}
\left\langle p_{\lambda},p_{\mu}\right\rangle =z_{\mu}\delta_{\lambda,\mu
}
\label{darij1.pf.t1.dualbases2}
\tag{9}
\end{equation}
for any $\lambda\in\operatorname{Par}$ and $\mu\in\operatorname{Par}$.




Now, fix a partition $\mu\in\operatorname{QPar}$. Then, Theorem 2 (applied to
$\lambda=\mu$) yields $z_{\mu}^{-1}p_{\mu}\in V\subseteq W$ (by Proposition
5). Hence, $z_{\mu}^{-1}p_{\mu}\in W=\Lambda_{\mathbb{Z}_{\left(  q\right)  }
}+J$. In other words, we can write $z_{\mu}^{-1}p_{\mu}$ in the form $z_{\mu
}^{-1}p_{\mu}=w_{1}+w_{2}$ for some $w_{1}\in\Lambda_{\mathbb{Z}_{\left(
q\right)  }}$ and some $w_{2}\in J$. Consider these $w_{1}$ and $w_{2}$.




Now, it is easy to see that $\left\langle f,w_{1}\right\rangle \in
\mathbb{Z}_{\left(  q\right)  }$. [\emph{Proof:} We have $f\in\Lambda
_{\mathbb{Z}_{\left(  q\right)  }}\cap\mathbb{Q}\left[  p_{i}\ \mid
\ i\not \equiv 0\mod q\right]  \subseteq\Lambda_{\mathbb{Z}
_{\left(  q\right)  }}$ and $w_{1}\in\Lambda_{\mathbb{Z}_{\left(  q\right)  }}$. Thus, $\left(f, w_1\right) \in \Lambda_{\mathbb{Z}_{\left(  q\right)  }}\times\Lambda_{\mathbb{Z}_{\left(  q\right)  }}$.
Hence, $\left\langle f,w_{1}\right\rangle \in\mathbb{Z}_{\left(  q\right)
}$, because the Hall inner product $\left\langle \cdot,\cdot\right\rangle $
sends $\Lambda_{\mathbb{Z}_{\left(  q\right)  }}\times\Lambda_{\mathbb{Z}
_{\left(  q\right)  }}$ to $\mathbb{Z}_{\left(  q\right)  }$.]




Furthermore, it is easy to see that $\left\langle f,w_{2}\right\rangle =0$.
[\emph{Proof:} We have $w_{2}\in J$; thus, we can write $w_{2}$ as a $\mathbb{Q}
$-linear combination of the family $\left(  p_{\lambda}\right)  _{\lambda
\in\operatorname{Par}\setminus\operatorname{QPar}}$ (since this family is a
basis of the $\mathbb{Q}$-vector space $J$). In other words, we can write
$w_{2}$ in the form
\begin{equation}
w_{2}=\sum_{\beta\in\operatorname{Par}\setminus\operatorname{QPar}}d_{\beta
}p_{\beta}
\label{darij1.pf.t1.w2=}
\tag{10}
\end{equation}
for some coefficients $d_{\beta}\in\mathbb{Q}$. Consider these $d_{\beta}$.
From \eqref{darij1.pf.t1.f=} and \eqref{darij1.pf.t1.w2=}, we obtain
\begin{align*}
\left\langle f,w_{2}\right\rangle  & =\left\langle \sum_{\lambda
\in\operatorname{QPar}}c_{\lambda}p_{\lambda},\sum_{\beta\in
\operatorname{Par}\setminus\operatorname{QPar}}d_{\beta}p_{\beta
}\right\rangle \\
& =\sum_{\lambda\in\operatorname{QPar}}\ \ \sum_{\beta\in\operatorname{Par}
\setminus\operatorname{QPar}}c_{\lambda}d_{\beta}\underbrace{\left\langle
p_{\lambda},p_{\beta}\right\rangle }_{\substack{=z_{\beta}\delta
_{\lambda,\beta}\\\text{(by \eqref{darij1.pf.t1.dualbases2})}}}\\
& =\sum_{\lambda\in\operatorname{QPar}}\ \ \sum_{\beta\in\operatorname{Par}
\setminus\operatorname{QPar}}c_{\lambda}d_{\beta}z_{\beta}\underbrace{\delta
_{\lambda,\beta}}_{\substack{=0\\\text{(since }\lambda\neq\beta
\\\text{(because }\lambda\in\operatorname{QPar}\\\text{whereas }\beta
\in\operatorname{Par}\setminus\operatorname{QPar}\text{))}}}\\
& =\sum_{\lambda\in\operatorname{QPar}}\ \ \sum_{\beta\in\operatorname{Par}
\setminus\operatorname{QPar}}c_{\lambda}d_{\beta}z_{\beta}0=0,
\end{align*}
qed.]




From $z_{\mu}^{-1}p_{\mu}=w_{1}+w_{2}$, we obtain
\begin{align*}
\left\langle f,z_{\mu}^{-1}p_{\mu}\right\rangle =\left\langle f,w_{1}
+w_{2}\right\rangle =\underbrace{\left\langle f,w_{1}\right\rangle }
_{\in\mathbb{Z}_{\left(  q\right)  }}+\underbrace{\left\langle f,w_{2}
\right\rangle }_{=0}\in\mathbb{Z}_{\left(  q\right)  }.
\end{align*}




On the other hand, from \eqref{darij1.pf.t1.f=}, we obtain
\begin{align*}
\left\langle f,z_{\mu}^{-1}p_{\mu}\right\rangle  & =\left\langle \sum
_{\lambda\in\operatorname{QPar}}c_{\lambda}p_{\lambda},z_{\mu}^{-1}p_{\mu
}\right\rangle =\sum_{\lambda\in\operatorname{QPar}}c_{\lambda}
\underbrace{\left\langle p_{\lambda},z_{\mu}^{-1}p_{\mu}\right\rangle
}_{\substack{=\delta_{\lambda,\mu}\\\text{(by \eqref{darij1.pf.t1.dualbases})}
}}=\sum_{\lambda\in\operatorname{QPar}}c_{\lambda}\delta_{\lambda,\mu}\\
& =c_{\mu}.
\end{align*}
Hence,
\begin{align*}
c_{\mu}=\left\langle f,z_{\mu}^{-1}p_{\mu}\right\rangle \in\mathbb{Z}_{\left(
q\right)  }.
\end{align*}




Forget that we fixed $\mu$. We thus have shown that $c_{\mu}\in\mathbb{Z}
_{\left(  q\right)  }$ for each $\mu\in\operatorname{QPar}$. In other words,
$c_{\lambda}\in\mathbb{Z}_{\left(  q\right)  }$ for each $\lambda
\in\operatorname{QPar}$. Hence, \eqref{darij1.pf.t1.f=} becomes
\begin{align*}
f=\sum_{\lambda\in\operatorname{QPar}}\underbrace{c_{\lambda}}_{\in
\mathbb{Z}_{\left(  q\right)  }}p_{\lambda}\in\sum_{\lambda\in
\operatorname{QPar}}\mathbb{Z}_{\left(  q\right)  }p_{\lambda}=\mathbb{Z}
_{\left(  q\right)  }\left[  p_{i}\ \mid\ i\not \equiv 0\mod 
q\right]
\end{align*}
(by the definition of $\operatorname{QPar}$).




Forget that we fixed $f$. We thus have shown that $f\in\mathbb{Z}_{\left(
q\right)  }\left[  p_{i}\ \mid\ i\not \equiv 0\mod q\right]  $
for each $f\in\Lambda_{\mathbb{Z}_{\left(  q\right)  }}\cap\mathbb{Q}\left[
p_{i}\ \mid\ i\not \equiv 0\mod q\right]  $. In other words,
\begin{align*}
\Lambda_{\mathbb{Z}_{\left(  q\right)  }}\cap\mathbb{Q}\left[  p_{i}
\ \mid\ i\not \equiv 0\mod q\right]  \subseteq \mathbb{Z}_{\left(  q\right)  }\left[ p_{i}\ \mid\ i\not \equiv 0\mod q\right] .
\end{align*}
As explained, this completes the proof of Theorem 1. $\blacksquare$



\section*{Sources}
Statement: R. P. Stanley, \emph{Supplementary Exercises for Chapter 7},
version 4 June 2023, exercise 21, physical page 4;
\url{https://math.mit.edu/~rstan/ec/ch7supp.pdf}.
Official solution locator: version 24 March 2023, physical page 6;
\url{https://math.mit.edu/~rstan/ec/ch7suppsol.pdf}.

Proof: Darij Grinberg, MathOverflow answer 403672 to question 403153,
created 11 September 2021, edited 30 September 2021.
\url{https://mathoverflow.net/a/403672}.
The API records the answer license as CC BY-SA 4.0:
\url{https://creativecommons.org/licenses/by-sa/4.0/}.
The original HTML/MathJax answer is retained alongside this editable typeset transcription.
The statement and answer have separate authors and source records.
\end{document}
